\clearpage
\SourceLeaf{357}\pdfbookmark[1]{Posthumous Papers and Publication}{parabolisch-veroeffentlichung}
\begin{center}
{\large POSTHUMOUS PAPERS AND PUBLICATION.\par}
\bigskip
\pdfbookmark[2]{II. On the Calculation of the True Anomaly in Parabolic Motion}{parabolisch-anomalie}
\textbf{[II. On the Calculation of the True Anomaly in Parabolic Motion.]}\par
\bigskip
\pdfbookmark[3]{1. New Table of Parabolic Motion}{tabula-nova}
\textbf{[1. New Table of Parabolic Motion.]}\par
\medskip
\textbf{Explanation of the New Table of Parabolic Motion.}
\end{center}
\medskip
\pdfbookmark[4]{1.}{tabula-nova-001}
\begin{center}1.\end{center}
\noindent\hspace*{\parindent}As in Arts. 18–20 of Theoria Motus, we shall denote by \(q\) the distance from the Sun at perihelion, by \(v\) the true anomaly, by \(t\) the time elapsed since passage through perihelion, and finally by [\(k\)] the constant \(3548''{,}188\) (Art. [6]). Since, in the parabolic motion of a comet whose mass we neglect, the angular velocity expressed in arcseconds is in general
\[
\frac{\mathrm{d}v}{\mathrm{d}t}=\frac{k\BareRoot{2}}{q^{\frac32}}\cos\tfrac12v^4,
\]
we shall set for brevity
\[
\frac{k\BareRoot{2}}{q^{\frac32}}=A,
\]
so that \(A\) represents the angular velocity at perihelion, which it will be convenient to add to the parabolic elements of any comet. Now
\[
\log k\BareRoot{2}=3{,}700\ 5216.
\]
\par\medskip
\pdfbookmark[4]{2.}{tabula-nova-002}
\begin{center}2.\end{center}
\noindent\hspace*{\parindent}The relation between the true anomaly and the time is expressed by the formula
\SourceLeaf{358}
\[
\operatorname{tang}\tfrac12v+\tfrac13\operatorname{tang}\tfrac12v^3
=\frac{At}{2\mathpunct{.}206265''}=\frac{\pi}{1296000''}At
\]
where \(\pi\) denotes half the circumference of a circle of unit radius. The time is thus obtained directly from the true anomaly. For the inverse calculation, \textsc{Cardano}'s rule, with the introduction of two auxiliary angles \(P\), \(Q\), supplies the following formulas:
\[
\operatorname{tang}P=\frac{4\mathpunct{.}206265''}{3At}=\frac{864000''}{\pi At}
\]
\[
\operatorname{tang}Q=\BareRoot[3]{\operatorname{tang}\tfrac12P}
\]
\[
\operatorname{tang}\tfrac12v=2\operatorname{cotang}2Q.
\]
\par\medskip
\pdfbookmark[4]{3.}{tabula-nova-003}
\begin{center}3.\end{center}
\noindent\hspace*{\parindent}But since the latter problem occurs very frequently, it would be tedious always to resort to these formulas or to indirect methods. Astronomers therefore generally use an auxiliary table for this purpose. The usual table gives, for the argument
\[
\frac{t}{q^{\frac32}}=\frac{At}{k\BareRoot{2}}
\]
the true anomaly. The celebrated \textsc{Burckhardt} adopted the logarithm of that quantity\footnote{Connaiss. des tems 1816.}. The argument of \textsc{Barker}'s table is the true anomaly, beside which the corresponding value of the quantity
\[
\frac{k\pi\BareRoot{2}\mathpunct{.}t}{17280''\mathpunct{.}q^{\frac32}}=\frac{\pi At}{17280''}
\]
is given. In the new table appended to this work (which had already been constructed when the table of the celebrated \textsc{Burckhardt} was published), we took the logarithm of the quantity \(At\) as the argument, proceeding in steps of one thousandth of a unit. Since in parabolic motion the value of \(At\) increases from \(0\) to infinity, the table evidently could not have been made complete at either end in this way. We therefore took the quantity \(At\) itself as the argument at the beginning of the table, proceeding by hundreds up to the value \(40000\), to which the true anomaly \(11^0\ 2'\ 32''\) corresponds. The logarithms then run from \(4{,}6\) to \(8{,}2\), to which limits the anomalies \(10^0\ 59'\ 26''\) and \(169^0\ 0'\ 15''\) correspond.
\par
\SourceLeaf{359}
Larger anomalies will very rarely be needed. Nevertheless, to make the table complete at both ends, we added a third part, in which we took as the argument the quantity
\[
\BareRoot[3]{\frac{B}{t}}\mathpunct{\raisebox{.45ex}{,}}
\]
where
\[
B=\frac{1296000^4}{3\pi^4A}=\frac{1296000^4}{3\pi^4k\BareRoot{2}}q^{\frac32}.
\]
\par
This argument decreases from \(40000\) to \(0\), to which the anomalies \(168^0\ 49'\ 9''\) and \(180^0\) correspond. To find the logarithm of the constant \(B\), we give the following logarithms
\[
\log\frac{1296000^4}{3\pi^4}=[22{,}461\ 8207]
\]
\[
\log\frac{1296000^4}{3\pi^4k\BareRoot{2}}=[18{,}761\ 2991.]
\]
\par
Thus the range of application of the first part of the table extends
\[
\begin{array}{rll}
\text{from}&t=0&\text{to }[t=]\ 7{,}971\,q^{\frac32}\\
\text{that of the second part from}&t=7{,}934\,q^{\frac32}&\text{to }t=31585\,q^{\frac32}\\
\text{that of the third part from}&t=30061\,q^{\frac32}&\text{to infinity.}
\end{array}
\]
\par
We also note that the entire table was originally calculated to hundred-thousandths of an arcsecond of anomaly, so that the second digit could always be relied upon.
\par\medskip
\pdfbookmark[4]{4.}{tabula-nova-004}
\begin{center}4.\end{center}
\noindent\hspace*{\parindent}It will be worth considering the beginning and end of our table separately and expanding the anomaly in a series proceeding in powers of the argument.
\par
Let us set
\[
\operatorname{tang}\tfrac12v=\theta
\]
\[
\frac{At}{206265''}=\alpha
\]
and then we shall have
\[
\theta=\tfrac12\alpha-\tfrac13\theta^3.
\]
From this we obtain, with the aid of \textsc{Lagrange}'s theorem, in general
\SourceLeaf{360}
\[
\begin{aligned}
\theta^n={}&\left(\frac{\alpha}{2}\right)^n-\frac{n}{3}\left(\frac{\alpha}{2}\right)^{n+2}
+\frac{n\mathpunct{.}n+5}{3\mathpunct{.}6}\left(\frac{\alpha}{2}\right)^{n+4}\\
&-\frac{n\mathpunct{.}+n+7\mathpunct{.}n+8}{3\mathpunct{.}6\mathpunct{.}9}\left(\frac{\alpha}{2}\right)^{n+6}\\
&+\frac{n\mathpunct{.}n+9\mathpunct{.}n+10\mathpunct{.}n+11}{3\mathpunct{.}6\mathpunct{.}9\mathpunct{.}12}\left(\frac{\alpha}{2}\right)^{n+8}-\text{etc.}
\end{aligned}
\]
Substituting this expression for the individual powers \(\theta^n\) in the series
\[
v=2\theta-\tfrac23\theta^3+\tfrac25\theta^5-\tfrac27\theta^7+\text{etc.}
\]
we have
\[
\begin{aligned}
v={}&\alpha-\tfrac13\cdot\tfrac14\left(1+\tfrac33\right)\alpha^3\\
&+\tfrac15\cdot\tfrac1{16}\left(1+\tfrac53+\tfrac{5\mathpunct{.}6}{3\mathpunct{.}6}\right)\alpha^5\\
&-\tfrac17\cdot\tfrac1{64}\left(1+\tfrac73+\tfrac{7\mathpunct{.}8}{3\mathpunct{.}6}+\tfrac{7\mathpunct{.}8\mathpunct{.}9}{3\mathpunct{.}6\mathpunct{.}9}\right)\alpha^7\\
&+\tfrac19\cdot\tfrac1{256}\Bigl(1+\tfrac93+\tfrac{9\mathpunct{.}10}{3\mathpunct{.}6}+\tfrac{9\mathpunct{.}10\mathpunct{.}11}{3\mathpunct{.}6\mathpunct{.}9}\\
&\hspace{5em}+\tfrac{9\mathpunct{.}10\mathpunct{.}11\mathpunct{.}12}{3\mathpunct{.}6\mathpunct{.}9\mathpunct{.}12}\Bigr)\alpha^9-\text{etc.,}
\end{aligned}
\]
where the pattern of the series is evident. On determining the numerical values of the coefficients we obtain
\[
\begin{aligned}
v={}&\alpha-\tfrac16\alpha^3+\tfrac{13}{240}\alpha^5-\tfrac{43}{2016}\alpha^7\\
&+\tfrac{191}{20736}\alpha^9-\tfrac{479}{114048}\alpha^{11}+\text{etc.}
\end{aligned}
\]
\par
Here \(v\) must evidently be regarded as an arc expressed in units of the radius. If it is wanted in arcseconds, the series must be multiplied by \(206265''\), so that it becomes
\[
v=At-\tfrac16\frac{(At)^3}{206265^2}+\tfrac{13}{240}\frac{(At)^5}{206265^4}-\text{etc.}
\]
Small anomalies therefore differ little from the corresponding arguments.
\par
Moreover, the expression given above for \(\theta^n\) also supplies an elegant series for the radius vector \(r\), for we obtain
\[
\begin{aligned}
r={}&q(1+\theta\theta)\\
={}&q\Bigl(1+\tfrac14\alpha\alpha-\tfrac13\cdot\frac{\alpha^4}{8}+\frac{7}{3\mathpunct{.}6}\frac{\alpha^6}{32}\\
&-\frac{9\mathpunct{.}10}{3\mathpunct{.}6\mathpunct{.}9}\frac{\alpha^8}{128}
+\frac{11\mathpunct{.}12\mathpunct{.}13}{3\mathpunct{.}6\mathpunct{.}9\mathpunct{.}12}\frac{\alpha^{10}}{512}\\
&-\frac{13\mathpunct{.}14\mathpunct{.}15\mathpunct{.}16}{3\mathpunct{.}6\mathpunct{.}9\mathpunct{.}12\mathpunct{.}15}\frac{\alpha^{12}}{2048}+\text{etc.}\Bigr).
\end{aligned}
\]
It can be shown that both series converge as long as \(\alpha<\tfrac43\), or equivalently \(v<61^0\ 35'\ 45''{,}4\).
\par
